\section{Body-hinge Tay--Whiteley theorem}
\label{sec:body-hinge}

This chapter proves the body-hinge (equivalently panel-hinge)
Tay--Whiteley theorem in $n$-space, in the existence-of-realization
form. It is the natural follow-on to the body-bar Tay theorem of
\cref{sec:body-bar} and reduces to it through a single combinatorial
device --- replacing each hinge by a bundle of $\delta - 1$ coincident
bars --- so the chapter adds no new linear algebra, only the
hinge-to-bar reduction and its bookkeeping.

\subsection{Body-hinge frameworks and the hinge-to-bar reduction}
\label{sec:body-hinge-reduction}

A body-hinge framework in $\mathbb{R}^{n}$ is a collection of
$n$-dimensional rigid bodies connected by \emph{hinges}, where a hinge
is an $(n - 2)$-dimensional affine subspace: a pin-joint in $2$-space,
a line-hinge in $3$-space, a plane-hinge in $4$-space, and so on
(Tay~\cite{tay1989}, Katoh--Tanigawa~\cite{katohTanigawa2011}). The
relative motion of two bodies sharing a hinge is constrained to a
rotation about it. As in \cref{sec:body-bar} we work with
$d = \mathrm{bodyBarDim}\,n = \binom{n+1}{2} = n(n+1)/2$, the dimension
of the space of infinitesimal screw motions of a single body in
$\mathbb{R}^{n}$; abbreviate $\delta = d$.

The key reduction (Tay~\cite{tay1989}~\S7, Whiteley~\cite{whiteley1988}):
a hinge constrains \emph{all but one} of the $\delta$ relative screw
freedoms of the two bodies it joins, so a single hinge behaves like a
bundle of $\delta - 1$ body-bars whose two-extensor coordinates lie in
the $(\delta-1)$-dimensional subspace dual to the hinge. Combinatorially,
this is captured by the \emph{parallel-edge multiplication} of the
underlying multigraph.

\begin{definition}[Edge multiplication $(\delta-1)\cdot G$]
  \label{def:edge-multiply}
  \lean{Graph.edgeMultiply, Graph.vertexSet_edgeMultiply, Graph.edgeMultiply_edgeSet_ncard, Graph.spanningVerts_edgeMultiply}
  \leanok
  \uses{def:graph-sparse}
  For a multigraph $G$ and $m \in \mathbb{N}$, let $m \cdot G$ be the
  multigraph on the same vertex set in which each edge of $G$ is
  replaced by $m$ parallel copies. In particular $V(m \cdot G) = V(G)$
  and $|E(m \cdot G)| = m\,|E(G)|$, and (for $m \ge 1$) a subset spans
  the same vertices in $m \cdot G$ as its image does in $G$. The edge
  type of $m \cdot G$ is $\beta \times \mathrm{Fin}\,m$; incidence is
  inherited from $G$ on the first coordinate.
\end{definition}

\begin{definition}[Body-hinge framework]
  \label{def:body-hinge-framework}
  \lean{Graph.BodyHingeFramework, Graph.BodyHingeFramework.toBodyBar,
    Graph.BodyHingeFramework.IsIndependent,
    Graph.BodyHingeFramework.IsInfinitesimallyRigid}
  \leanok
  \uses{def:edge-multiply, def:body-bar-framework}
  A \emph{body-hinge framework} in $\mathbb{R}^{n}$ on a multigraph
  $G$ assigns each hinge (edge) $e$ an $(n-2)$-dimensional affine
  subspace of $\mathbb{R}^{n}$, equivalently an
  $(\delta-1)$-dimensional space of bar coordinates dual to the hinge.
  Its rigidity map is the body-bar rigidity map
  (\cref{def:rigidity-map-body-bar}) of the induced body-bar framework
  on $(\delta - 1)\cdot G$ in which the $\delta - 1$ bars of each hinge
  carry a basis of that dual space. Infinitesimal rigidity and
  independence are inherited from the induced body-bar framework
  (\cref{def:infinitesimally-rigid-body-bar,def:independent-body-bar}).
\end{definition}

\begin{lemma}[Sparsity transports across edge multiplication]
  \label{lem:edge-multiply-sparse}
  \lean{Graph.BodyHingeFramework.edgeMultiply_isSparse_iff}
  \leanok
  \uses{def:edge-multiply, def:graph-sparse}
  For $\delta = \mathrm{bodyBarDim}\,n$, the multigraph $G$ is
  generically realizable as an independent (resp.\ isostatic)
  body-hinge framework in $\mathbb{R}^{n}$ exactly when
  $(\delta - 1)\cdot G$ is $(\delta, \delta)$-sparse (resp.\ tight).
  The count matches: $|E((\delta-1)\cdot G)| = (\delta - 1)|E(G)|$, and
  $(\delta,\delta)$-tightness of $(\delta-1)\cdot G$ reads
  $(\delta - 1)|E(G)| + \delta = \delta|V(G)|$, i.e.\
  $(\delta - 1)|E(G)| = \delta(|V(G)| - 1)$.
\end{lemma}
\begin{proof}
  \leanok
  \uses{def:body-hinge-framework, thm:tay-witness}
  A body-hinge framework on $G$ \emph{is} a body-bar framework on
  $(\delta - 1)\cdot G$ --- the placement data and the induced rigidity
  map are literally the same (\cref{def:body-hinge-framework}), so the
  two existentials are the same data (a bijection, formalized as the
  \texttt{exists\_toBodyBar\_iff} transport). Independence and
  isostaticity therefore read off node-for-node, and the claim is
  \cref{thm:tay-witness} applied to $(\delta - 1)\cdot G$. The generic
  hinge two-extensor assignment that realizes the $\delta - 1$
  coincident bars of each hinge as an independent bundle spanning the
  hinge's dual space is the standard-basis witness on the tree packing
  of $(\delta - 1)\cdot G$, carried verbatim from
  \cref{sec:body-bar-tay}.
\end{proof}

\subsection{Tay--Whiteley body-hinge theorem (existence-of-realization form)}
\label{sec:body-hinge-tay}

The chapter target. We ship the existence-of-realization form,
matching the scope of \cref{thm:tay-witness}; Whiteley's
``almost all realizations are rigid'' irreducible-variety lift
(Proposition~6 in~\cite{whiteley1988}) is not pursued here either
--- the standard-basis specialization on the multiplied graph is
again sufficient for the witness form.

\begin{theorem}[Body-hinge Tay--Whiteley theorem, independent form]
  \label{thm:body-hinge-tay}
  \lean{Graph.BodyHingeFramework.body_hinge_tay}
  \leanok
  \uses{def:body-hinge-framework, lem:edge-multiply-sparse,
    thm:tay-witness, thm:tutte-nash-williams, cor:k-spanning-trees}
  Let $n \ge 2$ and $\delta = \mathrm{bodyBarDim}\,n = n(n+1)/2$. A
  multigraph $G$ carries an independent body-hinge framework in
  $\mathbb{R}^{n}$ iff $(\delta - 1)\cdot G$ is
  $(\delta, \delta)$-sparse, and an isostatic one iff
  $(\delta - 1)\cdot G$ is $(\delta, \delta)$-tight --- equivalently, by
  \cref{thm:tutte-nash-williams,cor:k-spanning-trees}, iff
  $(\delta - 1)\cdot G$ is the edge-disjoint union of $\delta$ forests
  (resp.\ spanning trees). This is the Tay--Whiteley body-hinge theorem
  (Tay~\cite{tay1989}~\S7, Whiteley~\cite{whiteley1988}).
\end{theorem}
\begin{proof}
  \leanok
  \uses{lem:edge-multiply-sparse, thm:tay-witness}
  By \cref{lem:edge-multiply-sparse} a body-hinge framework on $G$ is
  independent (resp.\ isostatic) iff the induced body-bar framework on
  $(\delta - 1)\cdot G$ is; apply \cref{thm:tay-witness} to that
  framework. The tree-packing reformulations are
  \cref{thm:tutte-nash-williams,cor:k-spanning-trees} applied to
  $(\delta - 1)\cdot G$.
\end{proof}

The panel-hinge specialization --- bodies replaced by $(n-2)$-dimensional
panels, the $(n-1,2)$-frameworks of Tay~\cite{tay1989} --- shares the
same combinatorial characterization, so \cref{thm:body-hinge-tay}
covers both phrasings in the generic case. The further restriction in
which the hinges incident to each body are forced to be concurrent
(the \emph{panel-and-hinge} / \emph{molecular} model) is \emph{not}
covered here: it is the content of the molecular conjecture of Tay and
Whiteley, proved by Katoh--Tanigawa~\cite{katohTanigawa2011}, and is
the longer-horizon target of a future phase.
